\chapter{Answers to Selected Exercises}
\label{app:solutions}

This appendix gives the answers to a selection of the exercises in the book.
The answers are one solution, not the only solution; if your program gives
the right output and is readable, it is correct, even if it differs from
what is here. Please make your own attempt before you look at an answer.

\section*{Chapter \ref{ch:first}}

\solutionhead{Exercise 1.1}
\begin{salamcode}
func main:
    println "Sara"
    println "Ahmadi"
    println "Kashan"
end
\end{salamcode}

\solutionhead{Exercise 1.3}
\begin{salamcode}
func main:
    println "The sum of 125 and 375:", 125 + 375
    println "The product of 24 and 7:", 24 * 7
    println "Hours in a week:", 7 * 24
end
\end{salamcode}
\begin{salamout}
The sum of 125 and 375: 500
The product of 24 and 7: 168
Hours in a week: 168
\end{salamout}

\subsection*{Exercise 1.4}
Three mistakes: the colon after \code{func main} is missing, \code{println}
is misspelled on the second line, and the string on the third line has no
closing quotation mark.

\section*{Chapter \ref{ch:variables}}

\solutionhead{Exercise 2.1}
\begin{salamcode}
func main:
    length := 5
    width := 4
    height := 3
    println "Floor area:", length * width, "square metres"
    println "Volume of the room:", length * width * height, "cubic metres"
end
\end{salamcode}
\begin{salamout}
Floor area: 20 square metres
Volume of the room: 60 cubic metres
\end{salamout}

\subsection*{Exercise 2.3}
The output is \code{7 1}. Follow it through: \code{a} becomes 3 and
\code{b} becomes 6. Then \code{a} becomes 7. Finally \code{b} becomes
\code{7 - 6}, that is 1.

\subsection*{Exercise 2.4}
\code{score} is declared without \code{mut} but is changed on the next
line; it should be \code{mut score := 10}. \code{player name} is declared
but never used; either use it in a \code{println} or delete its line.

\section*{Chapter \ref{ch:types}}

\solutionhead{Exercise 3.3}
\begin{salamcode}
func main:
    total minutes := 135
    hours := total minutes / 60
    minutes := total minutes % 60
    println hours, "hours and", minutes, "minutes"
end
\end{salamcode}
\begin{salamout}
2 hours and 15 minutes
\end{salamout}

\section*{Chapter \ref{ch:operators}}

\solutionhead{Exercise 4.2}
\begin{salamcode}
func main:
    number := 4729
    ones := number % 10
    tens := (number / 10) % 10
    hundreds := (number / 100) % 10
    thousands := number / 1000
    println "Sum of the digits:", ones + tens + hundreds + thousands
end
\end{salamcode}
\begin{salamout}
Sum of the digits: 22
\end{salamout}

\section*{Chapter \ref{ch:conditions}}

\solutionhead{Exercise 5.3}
\begin{salamcode}
func main:
    hour := 15
    if hour < 12:
        println "Good morning"
    else hour < 17:
        println "Good afternoon"
    else hour < 20:
        println "Good evening"
    else:
        println "Good night"
    end
end
\end{salamcode}
\begin{salamout}
Good afternoon
\end{salamout}

\subsection*{Exercise 5.6}
Three lines are printed: ``greater than ten'', ``a multiple of five'' and
``a multiple of three''. The first chain ends with the first true condition,
so ``greater than five'' is not printed; but the two later \code{if}s are
separate and both are tested.

\section*{Chapter \ref{ch:loops}}

\solutionhead{Exercise 6.1}
\begin{salamcode}
func main:
    mut line := ""
    repeat 1 to 50 in n:
        if (n % 7) == 0:
            line = line + n + " "
        end
    end
    println line
end
\end{salamcode}
\begin{salamout}
7 14 21 28 35 42 49
\end{salamout}

\solutionhead{Exercise 6.2}
\begin{salamcode}
func main:
    mut product := 1
    repeat 1 to 10 in n:
        product *= n
    end
    println product
end
\end{salamcode}
\begin{salamout}
3628800
\end{salamout}

\subsection*{Exercise 6.6}
Six lines. The outer loop has three iterations, and in each of them the
inner loop prints only for \code{b} equal to 0 and 1, breaking at 2.

\section*{Chapter \ref{ch:functions}}

\solutionhead{Exercise 7.1}
\begin{salamcode}
func maximum(a: int, b: int): int:
    if a > b: ret a end
    ret b
end

func main:
    first two := maximum(12, 25)
    println maximum(first two, 19)
end
\end{salamcode}
\begin{salamout}
25
\end{salamout}
The largest of three numbers is the larger of ``the larger of the first
two'' and the third.

\subsection*{Exercise 7.6}
When \code{n} is zero, neither of the two conditions holds and the function
finishes without a \code{ret}; Salam gives a warning for this (appendix
\ref{app:errors}), and if you call the function with zero it returns
\code{null} instead of a number. The remedy: turn the second condition into a
plain \code{else:}, or put \code{ret n} at the end of the function.

\section*{Chapter \ref{ch:arrays}}

\solutionhead{Exercise 8.2}
\begin{salamcode}
func average(numbers: int[:]): f64:
    mut total := 0
    each n in numbers:
        total += n
    end
    count := numbers.len()
    ret (total as f64) / (count as f64)
end

func maximum(numbers: int[:]): int:
    mut largest := numbers[0]
    each n in numbers:
        if n > largest: largest = n end
    end
    ret largest
end

func main:
    grades := [17, 15, 19, 12, 20]
    println "Average:", average(grades[:])
    println "Maximum:", maximum(grades[:])
end
\end{salamcode}
\begin{salamout}
Average: 16.6
Maximum: 20
\end{salamout}

\subsection*{Exercise 8.6}
The output is \code{[10, 2, 30, 4, 50]}: the elements with even indices
(0, 2 and 4) are multiplied by ten.

\section*{Chapter \ref{ch:strings}}

\solutionhead{Exercise 9.4}
\begin{salamcode}
func main:
    text := "17,15,19,12"
    pieces := text.split(",")
    mut total := 0
    repeat pieces.len() in i:
        total += pieces.get(i).to_int()
    end
    count := pieces.len()
    println "Sum:", total
    average := (total as f64) / (count as f64)
    println "Average:", average
end
\end{salamcode}
\begin{salamout}
Sum: 63
Average: 15.75
\end{salamout}
